\documentclass[11pt,a4paper]{article}
\usepackage{gtu-el-yazisi}
\begin{document}
\notbilgi{Alan Teorisi}{Elektriksel dipol, kompleks sayılar, Gauss teoremi}{Ders7.pdf}

\sayfa{1}
\section{Dipol (Elektriksel)}

\[
p=q\cdot d\ \ (\mathrm{C\cdot m})=\text{sabit}
\qquad\qquad
\lim_{\substack{q\to\infty\\ d\to 0}} q\,d=p=\text{sabit}
\]

\sekilnotu[1]{Dipol geometrisi: $z$ ekseni yukarı, $y$ ekseni sağa, $x$ ekseni sol-aşağı. $z$ ekseni üzerinde $A(0,0,d)$ noktasında $+q$ yükü ve $B(0,0,-d)$ noktasında $-q$ yükü (orijinden her birine uzaklık $d$, süslü parantezle gösterilmiş; $B$ tarafında kesikli çizgi). Sağ üstte $P$ noktası; orijinden $P$'ye $r$, $A$'dan $P$'ye $r_1$, $B$'den $P$'ye $r_2$ doğruları. $r$ ile $+z$ ekseni arasındaki açı $\theta$, $r$ ile $-z$ yönü arasındaki açı $\pi-\theta$.}

\[
\lim_{d\to 0}\left(\frac{1}{r_1}-\frac{1}{r_2}\right)=\frac{1}{r}-\frac{1}{r}=0
\]

\[
V=\frac{q}{4\pi\varepsilon_0}\,\frac{1}{r_1}-\frac{q}{4\pi\varepsilon_0}\,\frac{1}{r_2}
\]

\[
\lim_{\substack{q\to\infty\\ d\to 0}}V=V_{dipol}=\frac{1}{4\pi\varepsilon_0}\lim_{\substack{q\to\infty\\ d\to\infty}}\frac{q\,d}{d}\left(\frac{1}{r_1}-\frac{1}{r_2}\right)\footnotemark
\]
\footnotetext{Claude notu: muhtemelen $d\to 0$ olmalı.}

\[
=\frac{1}{4\pi\varepsilon_0}\;\underline{\lim_{\substack{q\to\infty\\ d\to 0}}q\,d}\;\lim_{\substack{q\to\infty\\ d\to 0}}\frac{1}{d}\left(\frac{1}{r_1}-\frac{1}{r_2}\right)
\]

\[
V_{dipol}=\frac{p}{4\pi\varepsilon_0}\cdot\lim_{d\to 0}\frac{\left(\dfrac{1}{r_1}-\dfrac{1}{r_2}\right)}{d}
\]

\begin{align*}
\widehat{OAP}\quad r_1&=\sqrt{r^2+d^2-2rd\cos\theta}\\
\widehat{OBP}\quad r_2&=\sqrt{r^2+d^2-2rd\cos(\pi-\theta)}=\sqrt{r^2+d^2+2rd\cos\theta}
\end{align*}

\[
V_{dipol}=\frac{p}{4\pi\varepsilon_0}\lim_{d\to 0}\frac{\left(\dfrac{1}{r_1}-\dfrac{1}{r_2}\right)}{d}
\]
\[
=\frac{p}{4\pi\varepsilon_0}\lim_{d\to 0}\frac{\dfrac{1}{\sqrt{r^2+d^2-2rd\cos\theta}}-\dfrac{1}{\sqrt{r^2+d^2+2rd\cos\theta}}}{d}
\]
\[
=\frac{p}{4\pi\varepsilon_0}\lim_{d\to 0}\frac{d}{d(d)}\ \overset{0/0}{\Bigl[(r^2+d^2-2rd\cos\theta)^{-1/2}+(r^2+d^2+2rd\cos\theta)^{-1/2}\Bigr]}\footnotemark
\]
\footnotetext{Claude notu: muhtemelen iki terim arasında $-$ olmalı (bir sonraki satır buna uygun).}

\[
V_{dipol}=\frac{p}{4\pi\varepsilon_0}\lim_{d\to 0}\left[-\frac{1}{2}\,\frac{2d-2r\cos\theta}{(r^2+d^2-2rd\cos\theta)^{3/2}}+\frac{1}{2}\,\frac{2d+2r\cos\theta}{(r^2+d^2+2rd\cos\theta)^{3/2}}\right]
\]
\[
V_{dipol}=\frac{-p}{4\pi\varepsilon_0}\lim_{d\to 0}\left[\frac{d-r\cos\theta}{(r^2+d^2-2rd\cos\theta)^{3/2}}-\frac{d+r\cos\theta}{(r^2+d^2-2rd\cos\theta)^{3/2}}\right]\footnotemark
\qquad\kutula{\begin{aligned}2p&=2qd\\&=q(2d)\end{aligned}}
\]
\footnotetext{Claude notu: ikinci paydada muhtemelen $+2rd\cos\theta$ olmalı.}

\[
=-\frac{p}{4\pi\varepsilon_0}\left[-\frac{r\cos\theta}{r^3}-\frac{r\cos\theta}{r^3}\right]=\frac{p\cos\theta}{2\pi\varepsilon_0 r^2}
\qquad\qquad
V=\frac{2p\cos\theta}{4\pi\varepsilon_0 r^2}
\]
{\small $p\ \leftarrow$ dipolun moment yoğunluğu \qquad $2p$ (çemberlenmiş) $\leftarrow$ kitaptaki $p$}

\medskip
$z$ ekseni boyunca oluşturulmuş $p=qd$ moment yoğunluk bir noktasal dipolun potansiyeli

\sayfa{2}
\noindent$p=$ dipol moment yoğunluğu \qquad $[p]=\mathrm{C\,m}$

\noindent$[V_{dipol}]=\mathrm{V}$

\[
\kutula{\begin{array}{l}f(r,\theta,\phi)\\
\mathrm{grad}\,f=\dfrac{\partial f}{\partial r}\,\vec e_r+\dfrac{1}{r}\,\dfrac{\partial f}{\partial\theta}\,\vec e_\theta+\dfrac{1}{r\sin\theta}\,\dfrac{\partial f}{\partial\phi}\,\vec e_\phi\end{array}}
\]

\begin{multline*}
\vec E=-\mathrm{grad}\,V=-\frac{\partial V_{dipol}}{\partial r}\,\vec e_r-\frac{1}{r}\,\frac{\partial V_{dipol}}{\partial\theta}\,\vec e_\theta\\
-\frac{1}{r\sin\theta}\,\frac{\partial V_{dipol}}{\partial\phi}\,\vec e_\phi
\end{multline*}

\[
=\frac{2p\cos\theta}{4\pi\varepsilon_0 r^3}\,\vec e_r+\frac{1}{r}\,\frac{p}{4\pi\varepsilon_0 r^2}\sin\theta\,\vec e_\theta
\]
\textcolor{kirmizikalem}{kitaptaki $p$ $\downarrow$}
\[
=\frac{p}{4\pi\varepsilon_0}\Bigl[2\cos\theta\,\vec e_r+\sin\theta\,\vec e_\theta\Bigr]
\]

\textcolor{kirmizikalem}{Yada en baştan}
\[
\vec E=\frac{q}{4\pi\varepsilon_0}\,\frac{\vec e_{r_1}}{r_1^2}-\frac{q}{4\pi\varepsilon_0}\,\frac{1}{r_2^2}\,\vec e_{r_2}
\]
ama çok daha \okunamadi{zoooooor}!!

\formulkutu{\star\ \mathrm{rot}\,\vec E_{dipol}=0}

\section{Hadi kompleks sayılardan bahsedelim}

\noindent
\begin{minipage}[c]{0.46\linewidth}
\centering
\begin{tikzpicture}[line width=0.8pt]
  \draw[-{Stealth}] (-2.3,0) -- (2.6,0);
  \draw[-{Stealth}] (0,-1.1) -- (0,2.4) node[right] {$y$};
  \draw[-{Stealth}] (0,0) -- (1.3,1.0);
  \node[anchor=south west] at (1.3,0.95) {$z=x+iy$};
  \node at (0.42,0.78) {$r$};
  \draw (0.55,0) arc (0:38:0.55);
  \node at (0.82,0.18) {$\theta$};
  \draw (-0.08,1.0) -- (0.08,1.0);
  \node[anchor=east] at (-0.12,1.0) {$\mathrm{Im}(z)=y$};
  \node[anchor=north] at (1.3,-0.05) {$x=\mathrm{Re}(z)=r\cos\theta$};
\end{tikzpicture}
\end{minipage}\hfill
\begin{minipage}[c]{0.52\linewidth}
\begin{align*}
z&=\mathrm{Re}(z)+i\,\mathrm{Im}(z)\\
r&=\sqrt{x^2+y^2}=|z|\\
z&=r\cos\theta+i\,r\sin\theta\\
&=r\,[\cos\theta+i\sin\theta]\\
&=|z|\,[\cos(\arg z)+i\sin(\arg z)]
\end{align*}
\end{minipage}

\sayfa{3}
\section{Gauss Teoremi}

\[
\vec A=A_x\,\vec e_x+A_y\,\vec e_y+A_z\,\vec e_z
\]
\[
\underline{A_x,\ A_y,\ A_z\,(x,y,z)}
\]

\sekilnotu[3]{Düzensiz biçimli kapalı bir yüzey. Yüzey üzerindeki bir noktadan dışarı doğru $\vec n$ birim normali (yüzeye dik, dik açı işaretli); yanında $d\vec S=dS\cdot\vec n$. Yüzeyin sağ kenarının yanında ``$S$'' harfi; sağ kenardan sağa giden bir ok ``kapalı bir $S$ yüzeyi ve hacmi $v$'' yazısını gösteriyor (sonunda çok küçük bir ek var). Alttan bloğun içine doğru giden bir ok ``$v$ bölgesi'' yazısını gösteriyor.}

$v$\quad $S$ tarafından kuşatılan bölge

Gauss--Ostrogradski teoremi
\formulkutu{\int_{S}\vec A\cdot d\vec S=\int_{v}\mathrm{div}\,\vec A\cdot dv}

\[
\iint_{S}\vec A\cdot d\vec S=\iiint_{v}\mathrm{div}\,\vec A\cdot dv
\]

\end{document}
